\section{Related Work}
\label{sec:related-work}

The verified bibliography is organized by what each source is allowed to support. Gameplay-glitch
benchmarks frame the task and the temporal-glitch setting \cite{glitchbench,tempglitch}. Automated
game-testing work provides a broader software-engineering context in which testing goals and
oracles remain central \cite{politowski2022automatedtesting}. These sources motivate the problem;
they do not establish that the present repository solves production game QA.

The representation context comes from joint-embedding predictive architectures
\cite{ijepa,vjepa}, latent world models \cite{worldmodels,planet,dreamerv3}, and the external LeWM
paper \cite{lewm}. These sources explain why latent prediction and learned dynamics are relevant
background. They do not support any claim that this repository contributes a new architecture or
achieves control/planning results.

Video anomaly detection and OOD/calibration work are included to locate the evaluation problem
without importing their metrics as paper claims \cite{sultani2018ucfcrime,park2020mnad,
hendrycks2017baseline,guo2017calibration}. Leakage and reproducibility papers motivate the
paper's conservative split, threshold, and artifact-reporting discipline
\cite{kaufman2011leakage,kapoor2023leakage,pineau2021reproducibility}.

At this stage, the role of related work is not to claim a comprehensive literature survey. Instead,
it provides the minimum scholarly context necessary to position the paper without inventing
coverage. Table~\ref{tab:literature-matrix} records the current literature surface and the exact
claim boundary for each group.

\begin{table}[t]
\caption{Verified literature matrix. Sources provide context only unless a repository evidence
document separately supports a project-specific claim.}
\label{tab:literature-matrix}
\centering
\small
\begin{tabularx}{\textwidth}{@{}lX X@{}}
\toprule
Group & Supports & Does not support \\
\midrule
JEPA: \texttt{ijepa}, \texttt{vjepa} & Background on joint-embedding predictive architectures and video feature prediction. & Claims that this repository improves JEPA-family representation learning. \\
Latent world models: \texttt{worldmodels}, \texttt{planet}, \texttt{dreamerv3}, \texttt{lewm} & Context for learned latent dynamics and latent world-model framing. & Claims about control, planning, or broad world-model performance by this repository. \\
Video anomaly detection: \texttt{sultani2018ucfcrime}, \texttt{park2020mnad} & Background on video anomaly-detection formulations. & Direct comparability or superiority over surveillance-video anomaly systems. \\
Gameplay glitch detection: \texttt{glitchbench}, \texttt{tempglitch} & Motivation for gameplay glitch benchmarks and temporal glitch framing. & Any benchmark leaderboard, VLM, or broad glitch-detection superiority claim. \\
Automated game testing: \texttt{politowski2022automatedtesting} & Context that automated game testing remains workflow- and oracle-sensitive. & A production QA or human-replacement claim. \\
OOD and calibration: \texttt{hendrycks2017baseline}, \texttt{guo2017calibration} & Background for score-based detection and calibration caution. & Claims that LeWM surprise is calibrated probability or an OOD benchmark winner. \\
Leakage and reproducibility: \texttt{kaufman2011leakage}, \texttt{kapoor2023leakage}, \texttt{pineau2021reproducibility} & Rationale for split hygiene, learn-predict separation, and checklist-style reporting. & Claims that unaudited future experiments are already reproducible or leakage-free. \\
\bottomrule
\end{tabularx}
\end{table}


The literature section should remain cautious as the manuscript develops. In particular, this
scaffold does not argue that latent-surprise methods dominate alternative anomaly detectors; it
only notes that joint-embedding predictive architectures and latent world models are plausible
background for the repository's bounded evaluation.
